Our project has a constraint to fit within a very small space. We will utilize a partially
distributed integration strategy in order to mesh all of the components together. 
The primary electronics are an RFID module, a microcontroller unit (MCU), a power supply,
and a linear actuating solenoid. All components except the solenoid will be housed in the 
plastic grip for an AR-15. The plastic grip has the benefit of not interfering with the
RFID signal thus reducing the complexity of a faraday cage effect that could block any 
electromagnetic signal from being identified. 


The solenoid will be housed inside the actual receiver of the AR-15. This separate 
location allows for the direct blocking of the fire control group. The actuation will be 
triggered via fine gauge wire passing through one of the holes where the grip physically 
attaches to the receiver. 


Power management is one of the most challenging parts of this entire project. The RFID 
module can consume up to 60mA during its operation and the solenoid up to 200mA. To 
compensate for these power requirements we will employ an interrupt driven state machine. 
A low power timer will  periodically activate the RFID module to check for a tag, if one 
is detected it will activate hardware interrupt on the MCU to check if the tag is 
authorized. If the tag is in the authorized list the MCU will then actuate the solenoid 
to remove the block of the fire control group, ultimately allowing the firing of the AR-15.